\begin{abstract}
Contemporary artificial intelligence frameworks are predominantly optimized for specific hardware targets, creating tight couplings between algorithmic expression and physical execution. This coupling is antithetical to the vision of non-anthropocentric AI, which demands that knowledge genesis be substrate-agnostic and self-optimizing. We introduce APEIRON-Core, a multi-backend hardware abstraction layer (HAL) designed to provide a unified, low-overhead interface to four distinct computational paradigms: classical CPU/GPU, reconfigurable logic (FPGA), and quantum processing. The architecture is centered around a pluggable factory pattern with automatic backend selection, graceful fallback, and comprehensive observability. We detail the implementation of each backend—including a vectorized CPU backend with \(O(1)\) field caching, a PyTorch-based CUDA backend, a PYNQ-driven FPGA backend with fixed-point hardware integration, and a Qiskit-based quantum backend with corrected SWAP-test and partial trace operations. All backends are wrapped in a unified error-handling decorator that provides exponential backoff, circuit breaking, and Prometheus-compatible metrics. A novel energy-aware cost model links computational effort to thermodynamic expenditure, embedding an ethical resource constraint directly into the execution layer. We present an extensive experimental evaluation across all four backends, quantifying overhead, speedup, and energy efficiency. Results confirm that the abstraction overhead is <2\% on CPU/GPU for field sizes >\(10^4\), that FPGA acceleration yields up to \(7.8\times\) speedup on interference calculations, and that the quantum backend achieves >99.9\% fidelity on SWAP-test overlap measurements. The energy monitor successfully throttles operations when carbon intensity exceeds configurable thresholds, demonstrating the viability of thermodynamic self-regulation. APEIRON-Core forms the foundational execution substrate for the larger APEIRON framework, enabling higher cognitive layers to operate transparently across the full spectrum of available computational matter.
\end{abstract}